\subsection{Radial volumes on scalar Hopf manifolds}
\label{subsec:radial-gauduchon-lck-volumes}
\index{Gauduchon metric!radial LCK volumes}
\index{Lee form!prescribed radial period}

Let $n\geq2$, put $m=n-1$, and consider the scalar Hopf manifold and its projection to $\mathbb P^m$,
\[
 H_n=(\mathbb C^n\setminus\{0\})/\langle z\mapsto\mathrm e z\rangle,
 \qquad q:H_n\longrightarrow\mathbb P^m.
\]
Let $\omega_{\mathrm{FS}}$ be the Fubini--Study form on $\mathbb P^m$, normalized by $\int_{\mathbb P^m}\omega_{\mathrm{FS}}^m=1$. The function $r=\log|z|^2$ is well defined modulo $2$ on $H_n$, so its differential is a global one-form. Set
\[
 \alpha=q^*\omega_{\mathrm{FS}}=\ddc r,\qquad
 \beta=\mathrm dr\wedge\dc r.
\]
Both forms are defined on $H_n$. The form $\alpha$ has rank $m$, and its kernel is the complex radial line; the form $\beta$ has rank one, and its kernel is the complementary horizontal subspace. Thus $\alpha^m\wedge\beta$ is a positive volume form. On each elliptic fiber of $q$, where $z$ moves along a complex line through the origin, $\dc r=\mathrm d\arg z/(2\pi)$, so
\[
 \int_{q^{-1}(x)}\beta=2,\qquad
 \int_{H_n}\alpha^m\wedge\beta=2.
\]
We consider all metrics of the radial form
\begin{equation}\label{eq:radial-lck-family}
 \omega=C(r)\alpha+D(r)\beta,\qquad
 C,D\in C^\infty(\mathbb R/2\mathbb Z),\quad C,D>0.
\end{equation}
The circle $\mathbb R/2\mathbb Z$ of values of $r$ is oriented by increasing $r$. The \emph{Lee period} of such a metric is the integral of its Lee form over a radial loop along which $r$ increases from $0$ to $2$.

\Needspace{6\baselineskip}
\begin{lemma}\label{lem:radial-lck-gauduchon-equation}
Every metric of the form \eqref{eq:radial-lck-family} is strictly LCK. Its Lee form and Lee period are
\[
 \lambda=l(r)\,\mathrm dr,\qquad
 l=\frac{C'-D}{C},\qquad
 L=-\int_0^2\frac{D}{C}\,\mathrm dr<0.
\]
The metric is Gauduchon if and only if $C^{m-1}(C'-D)=k$ for a constant $k$. In terms of the positive smooth $2$-periodic function $F=C^mD$, this condition is equivalent to
\begin{equation}\label{eq:radial-lck-periodic-ode}
 C^mC'=kC+F,\qquad D=\frac{F}{C^m}.
\end{equation}
Every positive periodic solution has $k<0$, and
\begin{equation}\label{eq:radial-lck-period-identities}
 \begin{gathered}
 k\int_0^2C\,\mathrm dr+\int_0^2F\,\mathrm dr=0,
 \\
 l=kC^{-m},\qquad
 L(C)=-\int_0^2\frac{F}{C^{m+1}}\,\mathrm dr.
 \end{gathered}
\end{equation}
\end{lemma}
\begin{proof}
The differential identities are
\[
 \begin{gathered}
 \alpha^{m+1}=\beta^2=0,\qquad
 \mathrm d\beta=-\mathrm dr\wedge\alpha,\\
 \dc\beta=-\alpha\wedge\dc r,\qquad
 \ddc\beta=-\alpha^2.
 \end{gathered}
\]
They give $\mathrm d\omega=(C'-D)\mathrm dr\wedge\alpha=\lambda\wedge\omega$. The closed form $\lambda$ is the Lee form, since wedging one-forms with $\omega$ is injective. The integral of $C'/C$ over a period is zero, which proves the formula for $L$; since $L\neq0$, the Lee form is not exact. Two radial Lee forms have the same de Rham class exactly when their periods agree: a difference with zero period has a periodic primitive, and the radial loop detects the converse.

For a smooth function $G$ of $r$, the chain rule gives $\ddc G=G'\alpha+G''\beta$. Also $\dc(G\beta)=-G\alpha\wedge\dc r$, so
\[
 \ddc(G\beta)=-G'\alpha\wedge\beta-G\alpha^2.
\]
Expanding the powers of $\omega$ therefore gives
\[
 \begin{aligned}
 \omega^m&=C^m\alpha^m+mC^{m-1}D\alpha^{m-1}\wedge\beta,\\
 \omega^n&=nC^mD\alpha^m\wedge\beta,
 \end{aligned}
\]
and
\[
 \ddc(\omega^m)=
 \bigl[(C^m)''-m(C^{m-1}D)'\bigr]\alpha^m\wedge\beta.
\]
Since $(C^m)''-m(C^{m-1}D)'=m\bigl(C^{m-1}(C'-D)\bigr)'$, this vanishes exactly when $C^{m-1}(C'-D)$ is constant. Multiplying the identity $C^{m-1}(C'-D)=k$ by $C$ gives \eqref{eq:radial-lck-periodic-ode}. Since $C^mC'$ has zero integral over a period, integrating \eqref{eq:radial-lck-periodic-ode} gives the first identity in \eqref{eq:radial-lck-period-identities}, which forces $k<0$. The formulas for $l$ and $L$ follow from the formula for the Lee form.
\end{proof}

\Needspace{6\baselineskip}
\begin{lemma}\label{lem:radial-lck-periodic-order}
\index{Ordinary differential equation!periodic radial volumes}
For fixed $F>0$ and each $k<0$, there is at most one positive $2$-periodic solution of \eqref{eq:radial-lck-periodic-ode}. If $k_2>k_1$, the corresponding solutions satisfy $C_2>C_1$ everywhere and $L(C_2)>L(C_1)$.
\end{lemma}
\begin{proof}
Suppose $k_2>k_1$, and put $R=C_2-C_1$. At a zero of $R$, the equation gives $R'=(k_2-k_1)C_1^{1-m}>0$. Such zeros are isolated, so there are finitely many on the circle. A periodic differentiable function cannot cross zero only upwards, since returning around the circle would require a downward crossing. Hence $R$ has no zero and a constant strict sign. The identity $\int_0^2C_i\,\mathrm dr=-\int_0^2F\,\mathrm dr/k_i$ shows that this sign is positive, and the formula for $L(C)$ in \eqref{eq:radial-lck-period-identities} then gives $L(C_2)>L(C_1)$.

For equal $k$, uniqueness for the smooth scalar equation $C'=kC^{1-m}+FC^{-m}$ on $C>0$ shows that two solutions which meet coincide. Two solutions which never meet are strictly ordered, which contradicts the equality of their integrals. This proves uniqueness for fixed $k$.
\end{proof}

\Needspace{6\baselineskip}
\begin{theorem}\label{thm:radial-gauduchon-lck-volume}
For every positive smooth $2$-periodic function $F$ and every $L_0<0$, there is exactly one metric of the form \eqref{eq:radial-lck-family} which is Gauduchon, has Lee period $L_0$, and satisfies
\[
 \omega^n=nF(r)\alpha^m\wedge\beta.
\]
\end{theorem}
\begin{proof}
Fix $k<0$, put $c=-k$, and let $f_0=\min F>0$ and $f_1=\max F$. The constant functions $f_0/c$ and $f_1/c$ are a subsolution and a supersolution, respectively, of $C'=C^{-m}(F-cC)$. By comparison, the forward flow of this equation preserves the compact interval $I_k=[f_0/c,f_1/c]$ of positive numbers, and existence and uniqueness for smooth ordinary differential equations give a continuous time-two map $I_k\to I_k$. Its difference from the identity is nonnegative at the left endpoint and nonpositive at the right one, so the intermediate value theorem gives a fixed point; if $f_0=f_1$, the constant $C=f_0/c$ is the solution. Since the equation is periodic, the solution through a fixed point extends to a smooth positive periodic solution. It is unique by \cref{lem:radial-lck-periodic-order} and satisfies
\begin{equation}\label{eq:radial-lck-ode-barriers}
 \frac{f_0}{-k}\leq C_k(r)\leq\frac{f_1}{-k}.
\end{equation}
By \cref{lem:radial-lck-periodic-order}, the period $L(k)=-\int_0^2F/C_k^{m+1}\,\mathrm dr$ is strictly increasing in $k$. It is also continuous. Indeed, for $k$ in a compact neighborhood of any $k_0<0$, \eqref{eq:radial-lck-ode-barriers} gives uniform positive lower and finite upper bounds for $C_k$, and the equation then bounds all derivatives. Any sequence $k\to k_0$ therefore has smoothly convergent subsequences, and the equation and uniqueness identify every limit with $C_{k_0}$. Thus $C_k\to C_{k_0}$ uniformly, which gives continuity of the integral.

The same bounds imply $L(k)\to0$ as $k\uparrow0$ and $L(k)\to-\infty$ as $k\to-\infty$. Hence exactly one $k$ has period $L_0$. For this $k$, set $D=F/C_k^m$. The identities of \cref{lem:radial-lck-gauduchon-equation} show that $C_k\alpha+D\beta$ has the required properties. Conversely, every metric of the form \eqref{eq:radial-lck-family} with these properties gives a positive periodic solution of the same equation, so the ordering of the periods also proves uniqueness.
\end{proof}

\Needspace{6\baselineskip}
\begin{proposition}\label{prop:radial-lck-positive-twisted-potential}
\index{Twisted potential!positive radial construction}
Let $\omega$ be a metric as in \cref{thm:radial-gauduchon-lck-volume}, with Lee form $\lambda$. There is a positive smooth function $U$ on $H_n$ with $\omega=D_\lambda U$; in particular, the class of $\omega$ in $H^{1,1}_{\mathrm{BC},\lambda}(H_n,\mathbb R)$ is zero. If $\omega_1,\omega_2$ are two such metrics, with positive potentials $U_1,U_2$ and cohomologous Lee forms $\lambda_1$ and $\lambda_2=\lambda_1+\mathrm da$, then for every choice of the additive constant in $a$,
\[
 \mathrm e^{-a}\omega_2-\omega_1=D_{\lambda_1}(\mathrm e^{-a}U_2-U_1).
\]
In particular, among metrics of the form \eqref{eq:radial-lck-family}, a solution of \eqref{eq:lck-coupled-equation} is determined by its volume form and Lee class, even if the potential $u$ in \eqref{eq:lck-coupled-equation} is not assumed radial.
\end{proposition}
\begin{proof}
Write $l=kC^{-m}<0$ and $f(r)=\int_0^rl(s)\,\mathrm ds$ for real $r$. Then $f(r+2)=f(r)+L_0$, where $L_0<0$ is the Lee period of $\omega$. The solution of $U'-lU=C$ given by
\[
 \begin{aligned}
 U(0)&=\frac{\mathrm e^{L_0}}{1-\mathrm e^{L_0}}
       \int_0^2\mathrm e^{-f(s)}C(s)\,\mathrm ds,\\
 U(r)&=\mathrm e^{f(r)}
       \left(U(0)+\int_0^r\mathrm e^{-f(s)}C(s)\,\mathrm ds\right)
 \end{aligned}
\]
is positive on $[0,2]$ and satisfies $U(2)=U(0)$. Since the equation is periodic, all derivatives also match at the endpoints, so $U$ defines a positive smooth function on $H_n$. By \eqref{eq:lck-gauge-operators}, its twisted Hessian is
\[
 \begin{split}
 D_\lambda U
 &=(U'-lU)\alpha+
   \bigl[U''-2lU'+(l^2-l')U\bigr]\beta\\
 &=C\alpha+(C'-lC)\beta=\omega.
 \end{split}
\]
This proves that the twisted class is zero. The identity
$\mathrm e^{-a}D_{\lambda_2}U_2=D_{\lambda_1}(\mathrm e^{-a}U_2)$ of \cref{sec:lck-volume-equations} gives the formula for the difference. For two metrics of the form \eqref{eq:radial-lck-family}, equal volume forms give the same function $F$, and equal Lee classes give the same period, so \cref{thm:radial-gauduchon-lck-volume} applies. This uniqueness concerns the metric, not the potentials representing it.
\end{proof}

\Needspace{6\baselineskip}
\begin{proposition}\label{prop:radial-lck-non-vaisman-volume}
\index{Vaisman metric!nonconstant radial obstruction}
If $F$ is nonconstant, the metric of \cref{thm:radial-gauduchon-lck-volume} is not Vaisman.
\end{proposition}
\begin{proof}
The radial complex line is orthogonal to $\ker\mathrm dr\cap\ker\dc r$. On a radial complex line, use a logarithmic coordinate $w=x+\mathrm i y$, so that $r=2x+\text{constant}$ and $\beta=\pi^{-1}\mathrm dx\wedge\mathrm dy$. The Riemannian metric of $\omega$ restricts there to $D\pi^{-1}(\mathrm dx^2+\mathrm dy^2)$, whence
\[
 |\mathrm dr|_\omega^2=\frac{4\pi}{D},\qquad
 |\lambda|_\omega^2=4\pi k^2\frac{C^{-2m}}{D}.
\]
A parallel Lee form has constant norm. Since $k\ne0$, the Vaisman condition would therefore imply $D=Q C^{-2m}$ for a constant $Q>0$. At a maximum and at a minimum of $C$, where $C'=0$, the Gauduchon identity would then give $kC^{m+1}+Q=0$. This equation has only one positive root, so the maximum and the minimum of $C$ coincide. Thus $C$, $D$ and $F=C^mD$ would all be constant, a contradiction.
\end{proof}

\Needspace{6\baselineskip}
\begin{example}\label{ex:radial-gauduchon-non-vaisman}
\index{Hopf manifold!non-Vaisman radial metrics}
For
\[
 \begin{gathered}
 0<\varepsilon<\min\left\{\frac12,
                \frac{(3/2)^{1-m}}{4\pi}\right\},
 \\
 C=1+\varepsilon\cos(\pi r),\qquad
 D=C'+C^{1-m},
 \end{gathered}
\]
the metric \eqref{eq:radial-lck-family} is Gauduchon and non-Vaisman, with Lee form $\lambda=-C^{-m}\mathrm dr$ and density $F=C^mC'+C$. The probability volume $\mathrm dV=\frac12\alpha^m\wedge\beta$ gives $\omega^n=2nF\,\mathrm dV$ and $\int_{H_n}F\,\mathrm dV=1$.
\end{example}
\begin{proof}
Since $1/2<C<3/2$ and $m\geq1$, we have
$D\geq(3/2)^{1-m}-\pi\varepsilon>\frac34(3/2)^{1-m}>0$.
The Gauduchon identity holds with $k=-1$, giving the Lee form and the positive density $F=C^mD$. Its values $F(0)=1+\varepsilon$ and $F(1)=1-\varepsilon$ differ, so \cref{prop:radial-lck-non-vaisman-volume} applies. The integral of $C^mC'$ over a period is zero and the average of $C$ is one, so integration along the fibers of $q$ gives $\int_{H_n}F\,\mathrm dV=1$.
\end{proof}

The volume form in this example is not invariant under the flow of the Lee field of the Vaisman metric $\alpha+\beta$, which translates $r$, so \cref{thm:lck-vaisman-volume} does not apply to it. For the new metric, the Gauduchon condition states that $lF/D=lC^m=k$ is constant. Measured with the volume form $\omega^n$, the flux of its Lee field $\lambda^\sharp$ through each level set of $r$, oriented by increasing $r$, is $4\pi n\,lF/D=4\pi nk$; here $k=-1$. The radial classification answers \cref{q:lck-gauduchon-volume} within the family \eqref{eq:radial-lck-family}; uniqueness among arbitrary LCK metrics remains open.
